\uuid{uNgV}
\titre{Solution générale/raccord}
\auteur{hanani ; mzari}
\organisation{exo7}
\licence{CC BY-NC-SA 4.0}
\datecreate{2020-01-20}
\niveau{L1}
\module{Analyse}
\chapitre{Equation différentielle}
\sousChapitre{Résolution d'équation différentielle du premier ordre}
\difficulte{3}
\theme{Equations du premier ordre, solution générale, raccord}
\exo7qcmid{lille-2.equations-differentielles.019}

\contenu{
\begin{question}
\qclassification{225.01}
\qtags{motcle=solution générale/raccord}
On considère l'équation différentielle $(E)$ : $(1-x^2)y'-(1+x)y=1$. Parmi les affirmations suivantes, cocher celles qui sont vraies :
\begin{answers}
\good{La solution générale de $(E)$ sur $]-\infty ,-1[$, $]-1,1[$ ou $]1,+\infty[$ est :
$$\displaystyle y=\frac{k+\ln |1+x|}{1-x},\; k\in \Rr.$$}
\bad{L'équation $(E)$ admet une solution sur $\Rr$.}
\bad{L'équation $(E)$ admet une solution sur $]-\infty ,1[$.}
\good{L'équation $(E)$ admet une unique solution sur $]-1,+\infty [$.}
\end{answers}
\begin{explanations}
La solution générale de l'équation homogène sur l'un de ces intervalles  est $\displaystyle y=\frac{k}{1-x}$, $k\in \Rr$. La variation de la constante implique que $\displaystyle y_0=\frac{\ln |1+x|}{1-x}$ est une solution particulière. En posant $x=-1$ dans $(E)$, on obtient $0=1$ ce qui est absurde. Donc $(E)$ n'admet pas de solution ni sur $\Rr$ ni sur $]-\infty ,1[$.
\vskip2mm
\noindent Si $y$ est une solution sur $]-1,+\infty [$, on aura :
$$y(x)=\left\{\begin{array}{lll}\displaystyle \frac{k_1+\ln (1+x)}{1-x}&\mbox{si}&x\in ]-1,1[ \\ \\ \displaystyle \frac{k_2+\ln (1+x)}{1-x}&\mbox{si}&x\in ]1,+\infty[ \end{array}\right.$$
et $y(1)=-1/2$. La continuité en $1$ implique que $k_1=k_2=-\ln 2$. Une telle fonction est aussi dérivable en $1$. Ainsi $(E)$ admet une unique solution sur $]-1,+\infty [$.
\end{explanations}
\end{question}
}
